\section{Method}
\label{sec:method}

Given a question $q$ and schema view $\mathcal{S}$, a learned mapper proposes a typed computation
$z$.  The online path grounds and executes $z$, then releases a traced value, terminates without an
answer, or returns a structured failure for one further proposal.  An unsupported result or a clean
empty result terminates without another proposal.  Offline curation applies the same contract and a
hidden typed oracle before a candidate may become positive supervision.  The oracle is absent from
online prompts and student targets.  Figure \ref{fig:online-offline} shows both paths.

\begin{figure}[t]
  \centering
  \IfFileExists{figs/fig02_online_offline.pdf}{%
    \includegraphics[width=\linewidth]{fig02_online_offline}%
  }{%
    \fbox{\parbox[c][2.00in][c]{0.94\linewidth}{\centering\small
      \textbf{Figure placeholder.} The online path is question and schema $\rightarrow$ Step 1 intent
      $\rightarrow$ deterministic compiler or conditional Step 2 graph planner $\rightarrow$
      grounding $\rightarrow$ exhaustive execution $\rightarrow$ release or no answer.  Offline,
      the hidden oracle additionally gates direct SFT, repair targets, preference pairs, hard
      negatives, and abstention targets.  The oracle is removed before export.}}
  }
  \caption{Online inference and offline supervision construction.  The online runtime compiles,
  grounds, executes, and checks a proposal without access to an oracle.  Offline curation adds
  expected status, answer shape, and hidden oracle checks when applicable, then stores overlapping
  supervision views without exporting oracle values.  A repair reruns every deterministic stage.}
  \label{fig:online-offline}
\end{figure}

\subsection{Online proposal and checking}
\label{sec:m-runtime}

\paragraph{Structured understanding.}
The first model receives the question and a keyword-selected schema slice, but no database rows or
oracle.  It returns a JSON intent program containing an answer signature, typed steps, literal
constraints, role-bearing entity mentions, dependencies, and any unsupported or ambiguous reason.
The validator checks step references, return type, answer signature, and dependency order.

\paragraph{Direct compilation and conditional graph planning.}
A valid intent is compiled deterministically unless it triggers a fast-path veto.  The two current
vetoes prevent an unsupported cue from coexisting with an answering plan and prevent a request for
record count from compiling to an entity-set count.  If direct compilation fails or is vetoed, a
second model receives the question, the structured briefing, schema context, and the supported
capabilities.  It returns a typed variable graph.

Each graph variable has a kind, semantic role, constraints, and dependencies.  A relation binds a
field emitted by one variable into a downstream filter.  The return specification names the input
variable, operation, answer field, comparator or grouping key, and guards.  For example, a
buyer--supplier--buyer question is represented as
\[
 R_b \longrightarrow S_b=\operatorname{distinctSupplier}(R_b)
 \longrightarrow R[\operatorname{supplier}\in S_b]
 \longrightarrow \operatorname{distinctBuyer}(R)\setminus\{b\}.
\]
Before compilation, a consistency checker compares the graph with explicit years, CPV codes,
categories, titles, numeric literals, role cues, comparison direction, and the requested count
target.  In evaluation, the planner may draw two structural samples, but resampling occurs only
after a compile or consistency failure that can be detected without an oracle.

The reported implementation has one fidelity boundary after this check.  A late check may add a
missing year, CPV code, or category to a flat query specification after a graph plan already exists,
while execution still gives priority to that graph plan.  This path can detect a discrepancy but
does not guarantee that the added condition reaches execution.  We therefore audit retention of
question constraints separately and do not treat the late check as a semantic guarantee.

\paragraph{Hard grounding and preflight.}
Grounding maps aliases to allowed backend fields, checks value shapes, resolves buyer and supplier
mentions, verifies emitted and bound fields, forces exhaustive reduction, and inserts the additive
guard required by monetary sums.  Field grounding uses declared aliases and local character
$n$-gram candidates with type checks; entity resolution requires a unique acceptable match rather
than blindly selecting the first candidate.  A cycle, missing variable, incompatible dependency,
ambiguous entity, illegal field, or unsafe aggregate produces a structured failure.

\paragraph{Deterministic execution.}
Variables execute in dependency order over the complete contract-record universe.  Source variables
filter records, entity variables emit distinct values, dependent record variables consume those
values through set membership, and scalar variables count, sum, compare, rank, or select.  Counts
deduplicate contract identifiers, sums use decimal arithmetic, and ties use stable ordering.
Independent variables at the same dependency level may run concurrently without changing their
denotation.  Each variable trace records its grounded constraints, status, dependencies, output
size, and emitted field.

\paragraph{Evidence, release, and bounded repair.}
Preflight checks schema legality, conflicts, and exhaustive execution.  Postflight checks result
cardinality and population disclosure.  Checks for answer sanity and evidence then decide whether an
answer card may be returned.  The card contains the answer, exact population statistics, and a
bounded sample of matching contract identifiers.  The sample supports audit and replay but is not
the complete contributing population for a large aggregate.  Eligible compile, grounding, or
execution failures can trigger one replan conditioned on feedback.  Every repaired proposal is
compiled, grounded, executed, and checked again.  States with a clean empty result, ambiguity, or
unsupported content terminate without further answer search.  Training data collection permits two repair rounds;
the controlled runtime evaluation permits one.

\subsection{Offline supervision from checked executions}
\label{sec:m-training}

The teacher pipeline runs the full runtime on 9,267 training questions drawn from the benchmark
construction process described in Section \ref{sec:m-benchmark}.  The stored harvest summary names
\texttt{gpt-5.4-nano} for Step 1 and \texttt{grok-4-1-fast-non-reasoning} for Step 2.  The audited
runner reconstructs temperature zero and budgets of at most two structural plan samples and two
feedback repairs.  The historical artifact does not preserve a full command or prompt version, so
these decoding and budget values are reconstructed settings rather than independently stored call
metadata.  The exact prompt and call contract and its artifact boundary are given in Appendix
\ref{app:prompts}.  A
plan is a positive answerable target only when it passes compilation, grounding, execution, and release
checks, has the required answer shape, and matches the hidden oracle.  We write the online and
offline decisions separately.
\begin{equation}
\begin{aligned}
 R(z;q,\mathcal{S},K) &= G(z;q,\mathcal{S},K)\land V(z,K)
 \land D(\operatorname{Exec}(z,K),\sigma(z)), \\
 A(z;q,\mathcal{S},K,y,\tau) &= R(z;q,\mathcal{S},K)
 \land H(\operatorname{Exec}(z,K),\tau)
 \land M(\operatorname{Exec}(z,K),y).
\end{aligned}
\label{eq:acceptance}
\end{equation}
Here $G$ covers compilation and grounding, $V$ covers execution and release checks, and $D$ checks
the result against the signature $\sigma(z)$ declared by the proposal.  These terms define the
online decision $R$ without access to the answer.  Offline acceptance $A$ additionally compares the
result with the expected answer shape $\tau$ through $H$ and with the typed oracle $y$ through $M$.
A candidate that fails either offline check cannot become a positive target.  Wrong answers may be
retained as hard negatives, while successful repairs and accepted or rejected pairs are stored
separately.

Each training question contributes one primary trace, with overlapping repair and preference views
stored separately.  After family and bucket caps, training uses 2,787 plan rows, 1,679 repair rows,
and 5,091 Step 1 programs.  RSFT continues the SFT adapter on a stochastic Step 2 student harvest
while cloud Step 1 remains fixed.  Its stored summary records temperature 0.7, four plan samples,
and at most two feedback repairs.  This collection budget differs from the deterministic checkpoint
evaluation protocol in Section \ref{sec:experiments}.

The student never receives the oracle answer.  Direct plan supervised fine tuning learns
$p(z^+\mid q,u,\mathcal{S})$, where $u$ is the Step 1 briefing and $\mathcal{S}$ is the schema
contract.  Repair training learns $p(z^+\mid q,u,z^-,f,\mathcal{S})$ from a failed plan and
structured feedback.  Direct preference optimisation ranks accepted plans above rejected plans for
the same input \citep{rafailov2023dpo}.  These objectives explain how an oracle can select training
examples without becoming an answer shortcut at inference.

\subsection{Procurement data pipeline}
\label{sec:m-data}

\paragraph{Ingestion and snapshot selection.}
The loader streams compressed JSON Lines files and flattens OCDS releases, parties, awards, and
contracts.  Blank and malformed lines are logged and skipped.  A contracting process can be
published more than once.  For each OCID we retain the release with the latest parseable date and
store the set of source years in which that OCID appeared.  Extraction later accepts only this
selected $(\text{OCID},\text{release-id})$ pair, which prevents an older version from re-entering
the graph.

\paragraph{Entity resolution.}
Organisation resolution favours precision.  A recognised official scheme and identifier forms a
canonical node.  A platform-specific FTS identifier is linked to an official node only when all
observations with the same OCID, normalised name, and party role identify one candidate.  Unresolved
observations next pass through a government lookup, exact normalised name plus region, and finally
exact name when no region is available.  Residual clusters receive stable hashed identifiers, and
empty normalised names remain singletons.  Jaro--Winkler similarity of at least 0.92 produces a
review candidate but never an automatic merge.  A later union--find pass applies reviewed edges
only after checking cluster-level identifier conflicts.  Every raw identifier is retained in an
alias map.

\paragraph{Extraction, graph construction, and validation.}
The extractor creates tender, lot, award, bid-statistic, text-evidence, and document-metadata tables.
Amount selection follows award, linked contract, then tender precedence and preserves value,
currency, and source.  Organisation, contract, CPV, and evidence nodes are connected by
\textsc{buyer-of}, \textsc{supplier-of}, \textsc{categorised-by}, and \textsc{evidence-for} edges.
Exact lot evidence is preferred; an OCID-level link is used when no lot link is available.  Before
writing the graph, the build checks primary keys, edge endpoints, alias targets, JSON fields,
additivity policy, and award coverage.  Buyer, supplier, category, and evidence coverage are
reported separately rather than silently filling missing values.  Figure \ref{fig:data-graph}
summarises the representation path.

\begin{figure}[t]
  \centering
  \IfFileExists{figs/fig03_data_graph.pdf}{%
    \includegraphics[width=\linewidth]{fig03_data_graph}%
  }{%
    \fbox{\parbox[c][1.55in][c]{0.94\linewidth}{\centering\small
      \textbf{Figure placeholder.} (a) Yearly OCDS releases $\rightarrow$ latest dated
      $(\text{OCID},\text{release-id})$ snapshot. (b) Official identifiers $\rightarrow$
      conservative alias rules $\rightarrow$ reviewed safe merges or unresolved singletons.
      (c) Buyer and supplier organisations linked through contract-award nodes to CPV and evidence
      nodes, with source metadata retained.}}
  }
  \caption{Procurement data and representation pipeline.  Version selection precedes exact rich
  extraction; precision-first entity resolution preserves an alias audit; and the resulting
  Parquet-backed property graph retains category and contract-level provenance links.  Fuzzy
  similarity creates review candidates rather than automatic merges.}
  \label{fig:data-graph}
\end{figure}

\subsection{Executable benchmark construction}
\label{sec:m-benchmark}

Questions begin as deterministic \emph{answer specifications}, not as free-form model outputs.  A
specification contains literal constraints, an answer operation and field, a deduplication key,
expected status, and provenance.  Stage 1 executes each candidate on the production record backend
and a separate reference index.  It rejects conflicting constraints, inconsistent record IDs,
non-unique factoids, unstable answers, and unsafe sums.  Stage 2 renders an accepted specification
as a natural question.  A second gate checks identifier leakage, literal preservation, answer type,
operation, and visible constraints.  The naturaliser never writes the oracle.

The initial process retained 9,044 of 10,000 specifications after naturalisation.  Targeted builders
then added coverage-fixed, extended-operation, bridge, identity, and explicit non-answer examples.
The final 12,828 rows are assigned by whole plan group, so a plan identifier cannot occur in both
training and evaluation.  Curation also balances operation and difficulty buckets and diversifies
question surfaces.  The current artifact has no train--final plan-ID overlap, although two exact
question--answer strings and five deduplication groups cross that boundary; we report this in
Section \ref{sec:discussion}.

\subsection{Single call algebra planner}
\label{sec:m-compose}

PACS and WTQ use a separate planner that emits one recursive expression or an explicit abstention.
The algebra has seven types.  They are $\records$, $\values$, $\groups$, $\numbertype$,
$\valuetype$, $\booltype$, and $\ranking$.  Its operators support filtering, projection, grouping,
count, sum, extrema, comparison, set operations, and computed set membership.  Validation rejects unknown
fields, incompatible children, trees deeper than 16, and trees larger than 64 nodes.  The primary
PACS protocol uses one temperature-zero call, no repair, and no guided decoding.

Compose-v3 training examples are accepted only when a generated tree type-checks, the production
and secondary executors agree, and the result is non-degenerate.  Its current export has 12,414
training and 253 validation rows.  The WTQ port replaces the procurement record view with a
per table adapter.  A typed view is used for predicates and arithmetic, while a raw view preserves
the displayed answer strings.  Neither protocol invokes the full understander, graph planner,
evidence verdict, or feedback loop.
